\documentclass[10pt,a4paper]{amsart}
\usepackage{amsmath,amssymb,amsthm,graphicx}
% Retained minimal preamble for the Stacks Project chapter varieties.tex.
% Verbatim excerpt of the chapter preamble: Xy-pic setup (native preamble lines 14-18), the native theorem
% declarations with the mathematical macros (native preamble lines 154-217), and the project own no-op
% definitions of the reference/slogan/history display environments (native wrapper.tex lines 1-3) that the
% counted statement uses for its slogan. Copied mechanically; nothing added, removed or redefined.
% For commutative diagrams we use Xy-pic
\usepackage[all]{xy}

% We use 2cell for 2-commutative diagrams.
\xyoption{2cell}
% Theorem environments.
%
\theoremstyle{plain}
\newtheorem{theorem}[subsection]{Theorem}
\newtheorem{proposition}[subsection]{Proposition}
\newtheorem{lemma}[subsection]{Lemma}

\theoremstyle{definition}
\newtheorem{definition}[subsection]{Definition}
\newtheorem{example}[subsection]{Example}
\newtheorem{exercise}[subsection]{Exercise}
\newtheorem{situation}[subsection]{Situation}

\theoremstyle{remark}
\newtheorem{remark}[subsection]{Remark}
\newtheorem{remarks}[subsection]{Remarks}

\numberwithin{equation}{subsection}

% Macros
%
\def\lim{\mathop{\mathrm{lim}}\nolimits}
\def\colim{\mathop{\mathrm{colim}}\nolimits}
\def\Spec{\mathop{\mathrm{Spec}}}
\def\Hom{\mathop{\mathrm{Hom}}\nolimits}
\def\Ext{\mathop{\mathrm{Ext}}\nolimits}
\def\SheafHom{\mathop{\mathcal{H}\!\mathit{om}}\nolimits}
\def\SheafExt{\mathop{\mathcal{E}\!\mathit{xt}}\nolimits}
\def\Sch{\mathit{Sch}}
\def\Mor{\mathop{\mathrm{Mor}}\nolimits}
\def\Ob{\mathop{\mathrm{Ob}}\nolimits}
\def\Sh{\mathop{\mathit{Sh}}\nolimits}
\def\NL{\mathop{N\!L}\nolimits}
\def\CH{\mathop{\mathrm{CH}}\nolimits}
\def\proetale{{pro\text{-}\acute{e}tale}}
\def\etale{{\acute{e}tale}}
\def\QCoh{\mathit{QCoh}}
\def\Ker{\mathop{\mathrm{Ker}}}
\def\Im{\mathop{\mathrm{Im}}}
\def\Coker{\mathop{\mathrm{Coker}}}
\def\Coim{\mathop{\mathrm{Coim}}}

% Boxtimes
%
\DeclareMathSymbol{\boxtimes}{\mathbin}{AMSa}{"02}

%
% Macros for moduli stacks/spaces
%
\def\QCohstack{\mathcal{QC}\!\mathit{oh}}
\def\Cohstack{\mathcal{C}\!\mathit{oh}}
\def\Spacesstack{\mathcal{S}\!\mathit{paces}}
\def\Quotfunctor{\mathrm{Quot}}
\def\Hilbfunctor{\mathrm{Hilb}}
\def\Curvesstack{\mathcal{C}\!\mathit{urves}}
\def\Polarizedstack{\mathcal{P}\!\mathit{olarized}}
\def\Complexesstack{\mathcal{C}\!\mathit{omplexes}}
% \Pic is the operator that assigns to X its picard group, usage \Pic(X)
% \Picardstack_{X/B} denotes the Picard stack of X over B
% \Picardfunctor_{X/B} denotes the Picard functor of X over B
\def\Pic{\mathop{\mathrm{Pic}}\nolimits}
\def\Picardstack{\mathcal{P}\!\mathit{ic}}
\def\Picardfunctor{\mathrm{Pic}}
\def\Deformationcategory{\mathcal{D}\!\mathit{ef}}
\newenvironment{reference}{}{}
\newenvironment{slogan}{}{}
\newenvironment{history}{}{}
\usepackage{cleveref}
\crefname{originalreference}{original reference}{original references}
\Crefname{originalreference}{Original reference}{Original references}
\makeatletter
\newlabel{constructions-lemma-projective-space}{{01NE}{source}{}{sourcepointer-1}{}}
\newlabel{constructions-lemma-projective-space@cref}{{[originalreference][0][]01NE}{[1][0][]source}{}{}{}}
\makeatother
\crefname{originallemma}{original lemma}{original lemmas}
\Crefname{originallemma}{Original lemma}{Original lemmas}
\makeatletter
\newlabel{lemma-dimension-fibres-locally-algebraic}{{0B2L}{source}{}{sourcepointer-2}{}}
\newlabel{lemma-dimension-fibres-locally-algebraic@cref}{{[originallemma][0][]0B2L}{[1][0][]source}{}{}{}}
\makeatother
\crefname{originallemma}{original lemma}{original lemmas}
\Crefname{originallemma}{Original lemma}{Original lemmas}
\makeatletter
\newlabel{lemma-intersection-in-affine-space}{{0B2Q}{source}{}{sourcepointer-3}{}}
\newlabel{lemma-intersection-in-affine-space@cref}{{[originallemma][0][]0B2Q}{[1][0][]source}{}{}{}}
\makeatother
\title{Nonempty closed subsets of projective space whose dimensions sum to at least n intersect, with a dimension bound for the intersection}
\author{The Stacks Project Authors}
\date{stacks-project:varieties.tex:local-snapshot-20261011; GNU FDL 1.2 or later}
\begin{document}
\maketitle
\noindent\textit{Editorial note.} All delivered mathematical text is the author's own, copied verbatim from the pinned Stacks Project chapter file varieties.tex; no statement or proof step was rewritten, reordered or removed.\par\medskip
\noindent\textit{Editorial note.} The counted span cites four cross-chapter labels, each recorded in assembly.external\_references with its canonical tag as a navigation pointer.\par\medskip
\noindent\textit{Editorial note.} The author's parenthetical 'details omitted' concerns only the elementary verification that pi is surjective and smooth of relative dimension one with irreducible fibres; the main cone argument is written out.\par\medskip
\noindent\textit{Allowed standard prerequisites (editorial).}\begin{itemize}
\item The affine-space intersection bound. Stacks Project, Varieties, Lemma lemma-intersection-in-affine-space (tag 0A2T). Provides the dimension bound once the cones have been arranged to meet at the origin.
\item Dimensions of fibres and of preimages under a smooth morphism. Stacks Project, Varieties, Lemma lemma-dimension-fibres-locally-algebraic (tag 0B2L). Gives the dimension of the fibres of the restrictions of V and W used to show they contain every fibre they meet.
\end{itemize}\medskip
\ifdefined\hypertarget\hypertarget{sourcepointer-1}{}\fi
\noindent\textit{Source reference.} Original reference 01NE: sources/stacks-chow/tags.txt, source lines 2155--2155. Cross-chapter standard prerequisite cited by the author inside the counted span; it is used as a pointer and its own proof is not part of this item.\par\medskip
\ifdefined\hypertarget\hypertarget{sourcepointer-2}{}\fi
\noindent\textit{Source reference.} Original lemma 0B2L: sources/stacks-chow/tags.txt, source lines 13934--13934. Cross-chapter standard prerequisite cited by the author inside the counted span; it is used as a pointer and its own proof is not part of this item.\par\medskip
\ifdefined\hypertarget\hypertarget{sourcepointer-3}{}\fi
\noindent\textit{Source reference.} Original lemma 0B2Q: sources/stacks-chow/tags.txt, source lines 13938--13938. Cross-chapter standard prerequisite cited by the author inside the counted span; it is used as a pointer and its own proof is not part of this item.\par\medskip
\begin{lemma}
\label{lemma-intersection-in-projective-space}
Let $k$ be a field and $n \geq 0$. Let $X, Y \subset \mathbf{P}^n_k$
be nonempty closed subsets. If $\dim(X) = r$ and $\dim(Y) = s$ and
$r + s \geq n$, then $X \cap Y$ is nonempty and
$\dim(X \cap Y) \geq r + s - n$.
\end{lemma}

\begin{proof}
Write $\mathbf{A}^n = \Spec(k[x_0, \ldots, x_n])$ and
$\mathbf{P}^n = \text{Proj}(k[T_0, \ldots, T_n])$.
Consider the morphism
$\pi : \mathbf{A}^{n + 1} \setminus \{0\} \to \mathbf{P}^n$
which sends $(x_0, \ldots, x_n)$ to the point $[x_0 : \ldots : x_n]$.
More precisely, it is the morphism associated to the pair
$(\mathcal{O}_{\mathbf{A}^{n + 1} \setminus \{0\}}, (x_0, \ldots, x_n))$,
see Constructions, Lemma \ref{constructions-lemma-projective-space}.
Over the standard affine open $D_+(T_i)$ we get the morphism
associated to the ring map
$$
k\left[\frac{T_0}{T_i}, \ldots, \frac{T_n}{T_i}\right]
\longrightarrow
k\left[T_0, \ldots, T_n, \frac{1}{T_i}\right] \cong
k\left[\frac{T_0}{T_i}, \ldots, \frac{T_n}{T_i}\right]
\left[T_i, \frac{1}{T_i}\right]
$$
which is surjective and smooth of relative dimension $1$
with irreducible fibres (details omitted).
Hence $\pi^{-1}(X)$ and $\pi^{-1}(Y)$ are nonempty closed subsets of
dimension $r + 1$ and $s + 1$. Choose an irreducible component
$V \subset \pi^{-1}(X)$ of dimension $r + 1$ and an
irreducible component $W \subset \pi^{-1}(Y)$ of dimension $s + 1$.
Observe that this implies $V$ and $W$ contain every fibre of $\pi$
they meet (since $\pi$ has irreducible fibres of dimension $1$
and since Lemma \ref{lemma-dimension-fibres-locally-algebraic}
says the fibres of $V \to \pi(V)$ and $W \to \pi(W)$ have dimension $\geq 1$).
Let $\overline{V}$ and $\overline{W}$ be the closure of
$V$ and $W$ in $\mathbf{A}^{n + 1}$. Since $0 \in \mathbf{A}^{n + 1}$
is in the closure of every fibre of $\pi$ we see that
$0 \in \overline{V} \cap \overline{W}$. By
Lemma \ref{lemma-intersection-in-affine-space}
we have $\dim(\overline{V} \cap \overline{W}) \geq r + s - n + 1$.
Arguing as above using Lemma \ref{lemma-dimension-fibres-locally-algebraic}
again, we conclude that $\pi(V \cap W) \subset X \cap Y$
has dimension at least $r + s - n$ as desired.
\end{proof}

\end{document}
